\documentclass[12pt,a4paper]{report}
\usepackage[utf8]{inputenc}
\usepackage[T1]{fontenc}
\IfFileExists{frenchb.ldf}{\usepackage[french]{babel}}{\usepackage[english]{babel}}
\usepackage{lmodern}
\usepackage[margin=2.5cm]{geometry}
\usepackage{graphicx}
\usepackage[table]{xcolor}
\usepackage{multirow}
\usepackage{listings}
\usepackage{booktabs}
\usepackage{array}
\usepackage{float}
\usepackage{caption}
\usepackage{fancyhdr}
\usepackage{hyperref}
\usepackage{amsmath}

\hypersetup{colorlinks=true,linkcolor=blue!50!black,urlcolor=blue!50!black}
\pagestyle{fancy}\fancyhf{}
\fancyhead[L]{TP2 -- simulateur}\fancyhead[R]{\leftmark}\fancyfoot[C]{\thepage}
\setlength{\headheight}{15pt}

\lstdefinestyle{py}{
  language=Python, basicstyle=\ttfamily\footnotesize, keywordstyle=\color{blue!70!black}\bfseries,
  commentstyle=\color{green!40!black}\itshape, stringstyle=\color{orange!70!black},
  numbers=left, numberstyle=\tiny\color{gray}, breaklines=true, frame=single,
  backgroundcolor=\color{gray!6}, showstringspaces=false, tabsize=4, columns=fullflexible,
  literate={é}{{\'e}}1 {è}{{\`e}}1 {ê}{{\^e}}1 {à}{{\`a}}1 {â}{{\^a}}1 {î}{{\^i}}1 {ô}{{\^o}}1 {ù}{{\`u}}1 {û}{{\^u}}1 {ç}{{\c c}}1 {É}{{\'E}}1 {È}{{\`E}}1 {À}{{\`A}}1 {°}{{\textdegree}}1
}
\lstdefinestyle{shell}{basicstyle=\ttfamily\footnotesize, frame=single, backgroundcolor=\color{black!5},
  breaklines=true, columns=fullflexible, language={},
  literate={é}{{\'e}}1 {è}{{\`e}}1 {ê}{{\^e}}1 {à}{{\`a}}1 {â}{{\^a}}1 {î}{{\^i}}1 {ô}{{\^o}}1 {ù}{{\`u}}1 {û}{{\^u}}1 {ç}{{\c c}}1 {É}{{\'E}}1 {È}{{\`E}}1 {À}{{\`A}}1 {°}{{\textdegree}}1}
\lstset{style=py}

% Emplacement de capture : image si le fichier existe, sinon cadre vide.
% \capture{fichier.png}{Légende}{Commande qui produit la capture}
\newcommand{\capture}[3]{%
  \begin{figure}[H]\centering
  \IfFileExists{#1}{\includegraphics[width=\linewidth,height=0.6\textheight,keepaspectratio]{#1}}{%
    \fbox{\parbox[c][6cm][c]{0.92\linewidth}{\centering\color{gray}\textbf{[ CAPTURE D'ÉCRAN À INSÉRER ]}\\[4pt]
    \small Fichier attendu : \texttt{\detokenize{#1}}\\ Commande : \texttt{\detokenize{#3}}}}}%
  \caption{#2}\end{figure}}

\begin{document}

% ============================ PAGE DE GARDE ============================
\begin{titlepage}
\centering
{\large [Nom de l'établissement / Université]}\\[0.4cm]
{\large [Département / Filière]}\\[3cm]
{\Huge\bfseries Compte rendu du TP2\par}
\vspace{0.6cm}
{\LARGE Documentation et tests unitaires\\d'un package Python\par}
\vspace{0.6cm}
{\large Package \texttt{simulateur} : simulation de trafic routier\par}
\vspace{3cm}
\begin{tabular}{rl}
\textbf{Réalisé par :} & [Prénom NOM] \\[4pt]
\textbf{Encadré par :} & [Nom de l'enseignant] \\[4pt]
\textbf{Année universitaire :} & 2025--2026 \\
\end{tabular}
\vfill
{\large Novembre 2026}
\end{titlepage}

\tableofcontents
\listoffigures
\listoftables

% ============================ INTRODUCTION ============================
\chapter*{Introduction}
\addcontentsline{toc}{chapter}{Introduction}
Le package \texttt{simulateur} simule le déplacement de véhicules sur un réseau routier modélisé
par un graphe. Avant de le distribuer, ce TP demande de \textbf{documenter} le code, de rédiger un
\textbf{README}, de mettre en place une suite de \textbf{tests unitaires avec pytest}, puis de
\textbf{corriger les anomalies} que ces tests révèlent.

Ce compte rendu présente successivement :
\begin{enumerate}
  \item l'organisation du package et le travail de documentation (docstrings, README, site Sphinx) ;
  \item la conception des tests : cas nominaux, cas limites, cas invalides et cohérence ;
  \item l'exécution des tests avant et après correction, avec le décompte des tests réussis et échoués ;
  \item la liste des anomalies détectées et des corrections effectuées.
\end{enumerate}

% ============================ CHAPITRE 1 ============================
\chapter{Le package \texttt{simulateur}}
\section{Organisation}
Le code initial, réparti en plusieurs dossiers, a été réorganisé selon la structure demandée par
l'énoncé. Un module interne \texttt{\_\allowbreak{}validation.py} a été ajouté pour centraliser les contrôles de
paramètres (même type d'exception partout, sans duplication de code).

\begin{lstlisting}[style=shell,caption={Structure du projet}]
simulateur/
|-- simulateur/
|   |-- __init__.py        API publique du package
|   |-- _validation.py     contrôles de type et de valeur (interne)
|   |-- vehicule.py        classe Vehicule
|   |-- route.py           classe Route
|   `-- simulation.py      ReseauRoutier, Analyseur, Simulation
|-- tests/
|   |-- conftest.py        fixtures partagées
|   |-- test_vehicule.py
|   |-- test_route.py
|   `-- test_simulation.py
|-- docs/              documentation Sphinx (conf.py, *.rst)
|-- data/config_reseau.json
|-- exemples/exemple_simple.py
|-- README.md
`-- pyproject.toml
\end{lstlisting}

\capture{01_structure.png}{Arborescence réelle du projet}{python generer_captures.py}

\section{Rôle de chaque module}
\begin{table}[H]\centering
\begin{tabular}{@{}l p{4.2cm} p{6.3cm}@{}}
\toprule
\textbf{Module} & \textbf{Classes} & \textbf{Rôle} \\ \midrule
\texttt{vehicule.py} & \texttt{Vehicule} & Itinéraire, position, vitesse ; avance pendant un pas de temps et change de route \\
\texttt{route.py} & \texttt{Route} & Tronçon orienté : capacité, occupation, vitesse autorisée, accueil des véhicules \\
\texttt{simulation.py} & \texttt{ReseauRoutier}, \texttt{Analyseur}, \texttt{Simulation} & Graphe et Dijkstra ; statistiques et congestion ; boucle de simulation et lecture de la configuration \\
\texttt{\_\allowbreak{}validation.py} & (fonctions) & Validation de nombres, d'entiers et de textes \\ \bottomrule
\end{tabular}
\caption{Modules du package}
\end{table}

\section{Règles du simulateur}
Les tests s'appuient sur les règles suivantes.
\begin{itemize}
  \item Capacité d'une route : $\max\!\bigl(1,\ \text{longueur}\times\text{nb\_\allowbreak{}voies}/25\bigr)$ ; occupation $=$ véhicules / capacité.
  \item Vitesse autorisée : $v=v_{\text{limite}}\times\max(0{,}1,\ 1-\text{occupation})$.
  \item Vitesse réelle d'un véhicule : $\min(v_{\max},\,v)$ ; distance $=\dfrac{v}{3{,}6}\Delta t$.
  \item Itinéraire : chemin au temps de parcours théorique minimal (Dijkstra).
  \item Un véhicule se trouve toujours sur la route prévue par son itinéraire, à une position comprise entre $0$ et la longueur de la route.
\end{itemize}

% ============================ CHAPITRE 2 ============================
\chapter{Documentation}
\section{Docstrings}
Tous les éléments du package sont documentés : \textbf{49 modules, classes, fonctions et méthodes publics}
(vérifié automatiquement avec le module \texttt{ast}, aucun manquant). Le format retenu est le style Google :
\begin{itemize}
  \item \textbf{module} : objectif du fichier ;
  \item \textbf{classe} : rôle, responsabilités et liste des attributs avec leur type ;
  \item \textbf{fonction / méthode} : objectif, \texttt{Args} (types), \texttt{Returns}, \texttt{Raises}.
\end{itemize}
Les commentaires dans le code ont été conservés uniquement là où un point mérite explication
(par exemple la conversion km/h $\to$ m/s ou le report de la partie fractionnaire de véhicules à créer).

\begin{lstlisting}[caption={Exemple de docstring complète (\texttt{route.py})}]
def mettre_a_jour_vehicules(self, delta_t):
    """Fait avancer tous les véhicules de la route pendant ``delta_t`` secondes.

    La vitesse autorisée est calculée une seule fois, avec la charge du début de pas.

    Args:
        delta_t (float): Durée du pas en secondes, positive ou nulle.

    Returns:
        list[Vehicule]: Véhicules ayant atteint ou dépassé la fin de la route.

    Raises:
        TypeError: Si ``delta_t`` n'est pas un nombre.
        ValueError: Si ``delta_t`` est négatif.
    """
\end{lstlisting}

\capture{02_docstring_help.png}{Documentation lue avec \texttt{help()}}{python -c "from simulateur import Route; help(Route.vitesse_autorisee)"}
\capture{03_docstring_module.png}{Docstring de module affichée}{python -c "import simulateur.route as m; print(m.__doc__)"}

\section{Génération de la documentation avec Sphinx}
Les docstrings ne sont pas seulement lisibles dans le code : \textbf{Sphinx} les transforme en un site
HTML navigable (référence de l'API, index, recherche). La documentation est dans le dossier \texttt{docs/}.

\begin{table}[H]\centering
\begin{tabular}{@{}l p{10cm}@{}}
\toprule
\textbf{Fichier} & \textbf{Rôle} \\ \midrule
\texttt{docs/conf.py} & Configuration : chemin du package, extensions, thème, langue française \\
\texttt{docs/index.rst} & Page d'accueil et sommaire (\texttt{toctree}) \\
\texttt{docs/utilisation.rst} & Installation, exemples, erreurs levées \\
\texttt{docs/api.rst} & Référence de l'API, générée avec \texttt{automodule} \\
\texttt{docs/tests.rst} & Commandes de test \\
\texttt{docs/Makefile}, \texttt{make.bat} & Commande \texttt{make html} (Linux/macOS, Windows) \\ \bottomrule
\end{tabular}
\caption{Fichiers de la documentation Sphinx}
\end{table}

Trois extensions sont activées dans \texttt{conf.py} :
\begin{itemize}
  \item \texttt{sphinx.ext.autodoc} : extrait la documentation des docstrings ;
  \item \texttt{sphinx.ext.napoleon} : comprend le style Google (\texttt{Args}, \texttt{Returns}, \texttt{Raises}) utilisé dans le package ;
  \item \texttt{sphinx.ext.viewcode} : ajoute un lien vers le code source de chaque fonction.
\end{itemize}

\begin{lstlisting}[caption={Extrait de \texttt{docs/api.rst}}]
Module ``route``
----------------

.. automodule:: simulateur.route
   :members:
\end{lstlisting}

\begin{lstlisting}[style=shell,caption={Génération de la documentation}]
pip install -e ".[docs]"
cd docs
make html                 # Windows : make.bat html
# ouvrir docs/_build/html/index.html
\end{lstlisting}

\capture{11_sphinx_build.png}{Compilation de la documentation avec \texttt{sphinx-build}}{python -m sphinx -b html docs docs/_build/html}
\capture{12_sphinx_index.png}{Page d'accueil de la documentation HTML}{ouvrir docs/_build/html/index.html}
\capture{13_sphinx_api.png}{Référence de l'API générée depuis les docstrings}{ouvrir docs/_build/html/api.html}

\section{Fichier README.md}
Le \texttt{README.md} contient : la présentation du package, ses fonctionnalités, les prérequis
(Python $\ge 3.9$, aucune dépendance d'exécution), la procédure d'installation
(\texttt{pip install -e ".[dev]"}), un exemple d'utilisation, le format du fichier de configuration,
les exceptions levées et les commandes de test. Une section du README explique aussi comment générer la documentation Sphinx. Un fichier \texttt{pyproject.toml} décrit le package
et la configuration de pytest.

\begin{lstlisting}[style=shell,caption={Installation et lancement des tests}]
python -m venv .venv
source .venv/bin/activate
pip install -e ".[dev]"
pytest -v
\end{lstlisting}

\capture{09_exemple.png}{Exemple d'utilisation du package}{python exemples/exemple_simple.py}

% ============================ CHAPITRE 3 ============================
\chapter{Conception des tests unitaires}
\section{Méthode}
Les tests ont été écrits \emph{à partir des spécifications de l'énoncé} puis exécutés d'abord sur le
code d'origine. Chaque échec ainsi observé révèle donc une validation manquante ou un défaut réel,
corrigé ensuite dans le package. Un même fichier de tests sert avant et après correction, ce qui
permet de mesurer l'effet des corrections et de vérifier l'absence de régression.

\section{Les quatre types de tests}
\begin{table}[H]\centering\small
\begin{tabular}{@{}p{2.4cm} p{3.2cm} p{8.3cm}@{}}
\toprule
\textbf{Type} & \textbf{Module} & \textbf{Exemples de tests} \\ \midrule
\multirow{3}{*}{Cas nominal}
 & \texttt{vehicule} & création, conversion km/h $\to$ m/s (3 jeux de données), temps de parcours, changement de route \\
 & \texttt{route} & capacité et occupation, ajout/retrait d'un véhicule, véhicules arrivés en bout de route \\
 & \texttt{simulation} & Dijkstra, construction depuis la configuration, simulation de 30 tours, génération selon le débit \\ \midrule
\multirow{3}{*}{Cas limite}
 & \texttt{vehicule} & vitesse nulle, durée nulle, plafonnement par $v_{\max}$, itinéraire d'une seule route \\
 & \texttt{route} & route vide, pleine (40/40) et surchargée (60/40), capacité minimale de 1, plancher de vitesse à 10\,\% \\
 & \texttt{simulation} & 0 tour, débit nul, chemin vers soi-même, trajet impossible \\ \midrule
\multirow{3}{*}{Cas invalide}
 & \texttt{vehicule} & $v_{\max}\le 0$, itinéraire vide, départ négatif, identifiant vide, mauvais type, $\Delta t<0$ \\
 & \texttt{route} & longueur ou limite $\le 0$, voies $\in\{0,-1,1{,}5\}$, origine $=$ destination, nom vide, mauvais type \\
 & \texttt{simulation} & \texttt{n\_\allowbreak{}tours}$<0$, \texttt{delta\_\allowbreak{}t}$\le 0$, débit négatif, plage de vitesse incohérente, fichier absent, doublons, intersection inconnue \\ \midrule
\multirow{3}{*}{Cohérence}
 & \texttt{vehicule} & vitesse $\le\min(v_{\max}, v_{\text{auto}})$, position jamais décroissante \\
 & \texttt{route} & vitesse décroissante avec la charge, toujours entre 10\,\% et 100\,\% de la limite \\
 & \texttt{simulation} & conservation des véhicules, reproductibilité (graine), position dans $[0,\text{longueur}]$ à chaque tour, itinéraire respecté \\ \bottomrule
\end{tabular}
\caption{Organisation des tests par type et par module}
\end{table}

\section{Tests de la simulation}
Les quatre points demandés par l'énoncé sont couverts :
\begin{itemize}
  \item \textbf{exécution valide} : \texttt{test\_\allowbreak{}simulation\_\allowbreak{}nominale} (30 tours, historique de 30 entrées, temps simulé 1800 s) ;
  \item \textbf{paramètres invalides} : \texttt{n\_\allowbreak{}tours}, \texttt{delta\_\allowbreak{}t}, débit, plage de vitesse, fichier ;
  \item \textbf{règles de déplacement} : \texttt{test\_\allowbreak{}vehicules\_\allowbreak{}evoluent\_\allowbreak{}selon\_\allowbreak{}les\_\allowbreak{}regles\_\allowbreak{}de\_\allowbreak{}deplacement}
        (progression sur la route, vitesse $\le$ limite et $\le v_{\max}$) ;
  \item \textbf{contraintes du réseau} : \texttt{test\_\allowbreak{}contraintes\_\allowbreak{}du\_\allowbreak{}reseau\_\allowbreak{}respectees\_\allowbreak{}a\_\allowbreak{}chaque\_\allowbreak{}tour}
        (60 tours, position dans la route, itinéraire respecté, aucun doublon).
\end{itemize}

\begin{lstlisting}[caption={Fixtures et test paramétré (\texttt{tests/})}]
@pytest.fixture
def reseau_abc():
    r = ReseauRoutier()
    for nom, x in (("A", 0), ("B", 1), ("C", 2)):
        r.ajouter_intersection(nom, x, 0)
    r.ajouter_route(Route("A-B", "A", "B", 1000, 50))
    r.ajouter_route(Route("B-C", "B", "C", 1000, 50))
    return r

@pytest.mark.parametrize("nb_veh, attendue", [
    (0, 50.0), (10, 37.5), (20, 25.0), (40, 5.0), (80, 5.0)])
def test_vitesse_selon_occupation(nb_veh, attendue):
    r = Route("A-B", "A", "B", 1000, 50)
    remplir(r, nb_veh)
    assert r.vitesse_autorisee() == pytest.approx(attendue)

@pytest.mark.parametrize("longueur", [0, -100])
def test_longueur_non_positive_refusee(longueur):
    with pytest.raises(ValueError):
        Route("A-B", "A", "B", longueur, 50)
\end{lstlisting}

% ============================ CHAPITRE 4 ============================
\chapter{Exécution et analyse des résultats}
\section{Commandes}
\begin{lstlisting}[style=shell]
pytest -v                          # tous les tests
pytest tests/test_vehicule.py -v   # un module
pytest tests/test_route.py -v
pytest tests/test_simulation.py -v
\end{lstlisting}

\section{Avant correction (code d'origine)}
Les 95 tests ont été lancés sur le code d'origine réorganisé mais non modifié.
Un test ne se termine jamais (boucle infinie, anomalie S1) : il a dû être interrompu par un délai maximal.

\begin{table}[H]\centering
\begin{tabular}{@{}lrrrr@{}}
\toprule
\textbf{Module} & \textbf{Tests} & \textbf{Réussis} & \textbf{Échoués} & \textbf{Bloqués} \\ \midrule
\texttt{test\_\allowbreak{}vehicule.py} & 25 & 17 & 8 & 0 \\
\texttt{test\_\allowbreak{}route.py} & 28 & 17 & 11 & 0 \\
\texttt{test\_\allowbreak{}simulation.py} & 42 & 32 & 9 & 1 \\ \midrule
\textbf{Total} & \textbf{95} & \textbf{66} & \textbf{28} & \textbf{1} \\ \bottomrule
\end{tabular}
\caption{Résultats avant correction}
\end{table}

\capture{00_pytest_avant.png}{Exécution sur le code d'origine : 28 échecs, 66 réussites, 1 test bloqué}{pytest -q (code d'origine)}

\section{Après correction}
\begin{table}[H]\centering
\begin{tabular}{@{}lrrr@{}}
\toprule
\textbf{Module} & \textbf{Tests} & \textbf{Réussis} & \textbf{Échoués} \\ \midrule
\texttt{test\_\allowbreak{}vehicule.py} & 25 & 25 & 0 \\
\texttt{test\_\allowbreak{}route.py} & 28 & 28 & 0 \\
\texttt{test\_\allowbreak{}simulation.py} & 42 & 42 & 0 \\ \midrule
\textbf{Total} & \textbf{95} & \textbf{95} & \textbf{0} \\ \bottomrule
\end{tabular}
\caption{Résultats après correction}
\end{table}

\capture{04_pytest_complet.png}{\texttt{pytest -v} après correction : 95 tests réussis}{pytest -v}
\capture{05_pytest_vehicule.png}{Tests du module \texttt{vehicule}}{pytest tests/test_vehicule.py -v}
\capture{06_pytest_route.png}{Tests du module \texttt{route}}{pytest tests/test_route.py -v}
\capture{07_pytest_simulation.png}{Tests du module \texttt{simulation}}{pytest tests/test_simulation.py -v}
\capture{08_pytest_resume.png}{Résumé final}{pytest -q}

\section{Absence de régression}
Une simulation de 60 tours avec la configuration fournie donne exactement le même résumé qu'avant
les corrections (36,09\,km/h de vitesse moyenne, 2\,867 trajets terminés, 733 véhicules en route,
congestion sur \texttt{F-E} et \texttt{D-E}). Les corrections n'ont donc pas modifié le comportement
sur des données valides.

% ============================ CHAPITRE 5 ============================
\chapter{Anomalies détectées et corrections}
Les anomalies sont de deux natures : des \emph{validations absentes} (le code acceptait silencieusement des
valeurs absurdes) et des \emph{défauts de comportement}, plus graves (S1 et S5).

\section{Défauts de comportement}
\begin{table}[H]\centering\small
\begin{tabular}{@{}l >{\raggedright\arraybackslash}p{3.8cm} >{\raggedright\arraybackslash}p{4.0cm} >{\raggedright\arraybackslash}p{4.0cm}@{}}
\toprule
\textbf{N°} & \textbf{Où} & \textbf{Constat} & \textbf{Correction} \\ \midrule
S1 & \texttt{ReseauRoutier.\allowbreak{}generer\_\allowbreak{}vehicule} & \textbf{Boucle infinie} (\texttt{while True}) si aucun couple d'intersections n'est relié par un chemin & Nouvelle méthode \texttt{\_\allowbreak{}tirer\_\allowbreak{}trajet} : essais aléatoires bornés, puis parcours exhaustif des couples ; \texttt{ValueError} si aucun trajet \\
S5 & \texttt{ReseauRoutier.\allowbreak{}mettre\_\allowbreak{}a\_\allowbreak{}jour} & Avec un grand pas de temps, un véhicule pouvait rester à une position supérieure à la longueur de sa route (le dépassement couvrant plusieurs routes courtes n'était traité qu'une fois) & Boucle \texttt{while} : le véhicule traverse autant de routes que nécessaire dans le même pas \\ \bottomrule
\end{tabular}
\caption{Anomalies de comportement}
\end{table}

\begin{lstlisting}[caption={Correction de S5 : passage de plusieurs routes en un seul pas}]
for route, v in sortants:
    depassement = v.position - route.longueur
    while True:
        suivante = v.nom_route_suivante()
        if suivante is None:                       # fin de l'itinéraire
            route.retirer_vehicule(v)
            v.temps_arrivee = temps
            v.vitesse = 0.0
            self.vehicules_arrives.append(v)
            break
        nouvelle = self.routes[suivante]
        v.changer_de_route(nouvelle, depassement)
        if v.position < nouvelle.longueur:
            break                                  # reste sur cette route
        depassement = v.position - nouvelle.longueur
        route = nouvelle
\end{lstlisting}

\section{Validations manquantes}
\begin{table}[H]\centering\small
\begin{tabular}{@{}l >{\raggedright\arraybackslash}p{3.9cm} >{\raggedright\arraybackslash}p{4.4cm} >{\raggedright\arraybackslash}p{3.6cm}@{}}
\toprule
\textbf{N°} & \textbf{Où} & \textbf{Constat (avant)} & \textbf{Comportement (après)} \\ \midrule
V1 & \texttt{Vehicule()} & \texttt{vitesse\_\allowbreak{}max} $\le 0$ accepté (2 tests) & \texttt{ValueError} \\
V2 & \texttt{Vehicule()} & itinéraire vide accepté & \texttt{ValueError} \\
V3 & \texttt{Vehicule()} & \texttt{temps\_\allowbreak{}depart} négatif accepté & \texttt{ValueError} \\
V4 & \texttt{Vehicule()} & identifiant vide accepté & \texttt{ValueError} \\
V5 & \texttt{Vehicule()} & \texttt{vitesse\_\allowbreak{}max="rapide"} non détecté à la création & \texttt{TypeError} \\
V6 & \texttt{Vehicule.\allowbreak{}avancer} & \texttt{delta\_\allowbreak{}t} négatif : le véhicule \emph{reculait} & \texttt{ValueError} \\
V7 & \texttt{Vehicule.\allowbreak{}avancer} & vitesse autorisée négative acceptée & \texttt{ValueError} \\ \midrule
R1 & \texttt{Route()} & longueur $\le 0$ acceptée (2 tests) & \texttt{ValueError} \\
R2 & \texttt{Route()} & limite de vitesse $\le 0$ acceptée (2 tests) ; la division de Dijkstra pouvait planter & \texttt{ValueError} \\
R3 & \texttt{Route()} & \texttt{nb\_\allowbreak{}voies} $=0$, $-1$ ou $1{,}5$ accepté (3 tests) & \texttt{ValueError} / \texttt{TypeError} \\
R4 & \texttt{Route()} & origine $=$ destination acceptée & \texttt{ValueError} \\
R5 & \texttt{Route()} & nom vide accepté & \texttt{ValueError} \\
R6 & \texttt{Route()} & \texttt{longueur="longue"} : erreur obscure (\texttt{ValueError: could not convert...}) & \texttt{TypeError} explicite \\
R7 & \texttt{Route.\allowbreak{}mettre\_\allowbreak{}a\_\allowbreak{}jour\_\allowbreak{}vehicules} & \texttt{delta\_\allowbreak{}t} négatif accepté & \texttt{ValueError} \\ \midrule
S2 & \texttt{plus\_\allowbreak{}court\_\allowbreak{}chemin} & intersection inconnue : résultat silencieux (\texttt{None} ou \texttt{[]}) & \texttt{ValueError} \\
S3 & \texttt{ajouter\_\allowbreak{}route} & route en double : l'ancienne était écrasée avec ses véhicules & \texttt{ValueError} \\
S4 & \texttt{ajouter\_\allowbreak{}intersection} & intersection en double écrasée & \texttt{ValueError} \\
S6 & \texttt{lancer\_\allowbreak{}simulation} & \texttt{n\_\allowbreak{}tours} négatif : 0 tour exécuté sans erreur & \texttt{ValueError} \\
S7 & \texttt{lancer\_\allowbreak{}simulation} & \texttt{delta\_\allowbreak{}t} $\le 0$ accepté (2 tests) & \texttt{ValueError} \\
S8 & \texttt{Simulation()} & débit \texttt{vehicules\_\allowbreak{}par\_\allowbreak{}minute} négatif accepté & \texttt{ValueError} \\
S9 & \texttt{Simulation()} & plage de vitesse \texttt{[130, 90]} acceptée & \texttt{ValueError} \\ \bottomrule
\end{tabular}
\caption{Validations ajoutées}
\end{table}

\paragraph{Bilan.} Ces anomalies correspondent aux \textbf{28 échecs} et au \textbf{test bloqué} du
tableau avant correction (V : 8 échecs, R : 11, S : 9 échecs et 1 blocage). Toutes les
validations passent par les trois fonctions de \texttt{\_\allowbreak{}validation.py}.

\begin{lstlisting}[caption={Exemple de validation (\texttt{route.py})}]
self.longueur = verifier_nombre("longueur", longueur, strictement_positif=True)
self.limite_vitesse = verifier_nombre("limite_vitesse", limite_vitesse,
                                      strictement_positif=True)
self.nb_voies = verifier_entier("nb_voies", nb_voies, minimum=1)
\end{lstlisting}

\capture{10_erreur_validation.png}{Message d'erreur clair après correction (\texttt{Route} avec longueur négative)}{python -c "from simulateur import Route; Route('A-B','A','B',-100,50)"}

\section{Choix de conception}
\begin{itemize}
  \item \textbf{Types} : un mauvais type lève \texttt{TypeError} ; une valeur hors domaine lève \texttt{ValueError}.
        Les booléens sont refusés comme nombres (sinon \texttt{True} passerait pour $1$).
  \item Un test écrit trop strictement (\texttt{nb\_\allowbreak{}voies=1.5} attendu en \texttt{ValueError}) a été corrigé : un
        décimal est un mauvais \emph{type}, donc \texttt{TypeError}. Le code était correct, la spécification du test ne l'était pas.
  \item Les paramètres \texttt{n\_\allowbreak{}tours} décimal ou texte levaient déjà une \texttt{TypeError} (par \texttt{range}) :
        ce n'était pas une anomalie.
\end{itemize}

% ============================ CONCLUSION ============================
\chapter*{Conclusion}
\addcontentsline{toc}{chapter}{Conclusion}
Le package \texttt{simulateur} est désormais entièrement documenté (docstrings complètes, README, \texttt{pyproject.toml}) et
validé par \textbf{95 tests pytest, tous réussis}. Écrire les tests à partir des spécifications avant de corriger
a mis en évidence \textbf{29 problèmes} sur le code d'origine : 28 tests en échec et un test bloqué, soit
deux défauts réels (une boucle infinie et un dépassement de route) et des validations absentes sur les
paramètres. Le comportement sur des données valides est inchangé.

\paragraph{Pistes d'amélioration.}
mesurer la couverture avec \texttt{pytest-cov}, ajouter une intégration continue, et tester
l'affichage graphique et l'export, hors du périmètre de ce TP.

\end{document}
